\documentclass[11pt, a4paper]{article}

\usepackage{amsmath}
\usepackage{amssymb}
\usepackage{cite}
\usepackage{hyperref}
\usepackage{geometry}

\geometry{margin=2.5cm}

\title{Encryption Scheme}
\author{}
\date{\today}

\begin{document}

\maketitle

\section{Yao's Garbled Circuits}

\section{Fairplay} \cite{fairplay}

\begin{enumerate}
    \item Encryption:
        \begin{enumerate}
            \item Generate pseudo-random pads: 
            $\text{SHA1}(w_i\parallel k\parallel x\parallel y)$ and 
            $\text{SHA1}(w_j\parallel k\parallel x\parallel y)$, 
            $w_i, w_j$ input wire labels, $k$ gate index, $x, y$ input bits. 
            Each pad is a 160-bit hash value.
            \item Apply XOR: $ciphertext=output\_label\oplus pad_i\oplus pad_j$.
        \end{enumerate}
    \item Decryption: $output\_label=ciphertext\oplus pad_i\oplus pad_j$.
    \item Pros:
        \begin{enumerate}
            \item SHA-1 supports variable-length inputs (AES uses block ciphers).
            \item Wire labels can be used to generate pads, without introducing
            additional keys.
            \item XOR is simple and efficient.
        \end{enumerate}
    \item Cons:
        \begin{enumerate}
            \item SHA-1 is outdated.
            \item Existing weakness: when $w_i=w_j \rightarrow pad_i=pad_j$,
            $pad_i\oplus pad_j=0$. \cite{duplicate_input_wires}
            \item Uses SHA-1 as a pseudo-random function.
        \end{enumerate}
\end{enumerate}

\section{Proof of Security} \cite{proof_of_security}

\begin{enumerate}
    \item Encryption:
        \begin{enumerate}
            \item Double encryption: 
            $C_{x,y} = E_{k_1^x}\left(E_{k_2^y}(k_3^{g(x,y)})\right)$, $k_1^x, k_2^y$ input wire labels, $k_3^{g(x,y)}$ output label, $E$ symmetric encryption.
            \item $E_k(x)=\langle r,\ f_k(r)\oplus x0^n\rangle$, $r$ random string (n-bits), $x$ plain text (n-bits).
        \end{enumerate}
    \item Note: requirements for secure encryption scheme: CPA security, secure 
    for multiple messages, elusive range, efficiently verifiable range, and double-encryption.
security.
\end{enumerate}




\bibliographystyle{plain}
\bibliography{references}

\end{document}